\section{Validation}

The system has been validated through a combination of automated testing and API-level verification, with the objective of verifying both the internal correctness of services and the proper behavior of exposed interfaces.

Two main testing approaches have been adopted: unit and integration testing of service logic, and API-level testing through interactive documentation.

\subsection{Automated Testing and Coverage}

Each service has been tested through automated tests designed to cover:
\begin{itemize}
  \item domain logic and service-layer behavior;
  \item validation rules and edge cases.
\end{itemize}

Integration tests of data persistence operations have been performed using a MongoDB in-memory implementation or by a MongoDB container automatically started (and torn down) during Gradle test execution.

To assess the effectiveness of the test suite, code coverage has been measured. The project achieves an overall coverage of at least 70\%, ensuring that the majority of the codebase is exercised during testing. This level of coverage provides reasonable confidence in the correctness of the implemented logic.

\subsection{API Testing}

In addition to automated tests, the system has been validated through direct interaction with the exposed APIs.

Each service provides a formal API specification compliant with the OpenAPI standard, which allows endpoints to be explored and tested interactively. This interface has been used to manually verify the correctness of request handling, response formats and error management.

API testing has been particularly useful for:
\begin{itemize}
  \item validating request/response workflows;
  \item testing edge cases and invalid inputs;
  \item verifying correctness of response formats.
\end{itemize}

The full set of APIs is available at: \url{https://eventonight.github.io/EvenToNight/openAPI/}.

Moreover, APIs have been tested in integration with the frontend, ensuring correct implementation and usage.
